\begin{question}[index=17,points=10]
问是否存在常数 $k$，使得 $x\to0$ 时，$\displaystyle\int_0^{2x}|t-x|\sin t\dd t$ 与 $k|x|^3$ 是等价无穷小？若存在，求 $k$ 的值；若不存在，请说明理由.
\end{question}
\QuestionBreak\WorkSpace

\begin{question}[index=18,points=12]
设球面 $\Sigma:x^2+y^2+z^2=2x$，计算 $\displaystyle I=\iint_\Sigma(2x^2+y^2+3z^2)\dd S$.
\end{question}
\QuestionBreak\WorkSpace\WorkPageBreak

\begin{question}[index=19,points=12]
设 $\displaystyle f(x)=\sum_{n=0}^\infty(-1)^na_nx^{2n+1}$ 满足 $f''(x)+f(x)=-2\sin x$，$f'(0)=2$.
\par（Ⅰ）求 $f(x)$；
\par（Ⅱ）求 $a_n$ 及级数 $\displaystyle\sum_{n=0}^\infty\frac{(-1)^na_n}{2^n}$ 之和.
\end{question}
\QuestionBreak\WorkSpace

\begin{question}[index=20,points=12]
设 $f(x)$ 在 $(-\infty,+\infty)$ 内二阶可导，且对任意 $h$，有 $f(x+h)+f(x-h)-2f(x)\ge0$.
\par（Ⅰ）证明：$f''(x)\ge0$；
\par（Ⅱ）证明：$\displaystyle\int_0^1f(x^n)\dd x\ge f\!\left(\frac1{n+1}\right)$，其中 $n$ 为正整数.
\end{question}
\QuestionBreak\WorkSpace\WorkPageBreak

\begin{question}[index=21,points=12]
设二次型 $f(x_1,x_2,x_3)=\vect X^{\mathrm T}\vect A\vect X=2x_1^2+2x_2^2+ax_3^2+2x_1x_2+2x_1x_3+2x_2x_3$，$\vect A^{\mathrm T}=\vect A$，$r(\vect A)=3$．向量 $\vect\alpha=(1,b,1)^{\mathrm T}$ 满足 $(\lambda\vect E-\vect A^*)\vect\alpha=\vect0$，其中 $b\ne1$，$\vect A^*$ 是 $\vect A$ 的伴随矩阵.
\par（Ⅰ）求 $a,b,\lambda$ 的值；
\par（Ⅱ）求正交变换 $\vect X=\vect Q\vect Y$，将 $g(x_1,x_2,x_3)=\vect X^{\mathrm T}(\vect A-\vect A^*)\vect X$ 化为标准形.
\end{question}
\QuestionBreak\WorkSpace

\begin{question}[index=22,points=12]
设随机变量 $X$ 与 $Y$ 相互独立，$X$ 服从参数为 $\lambda$（$\lambda>0$）的指数分布，$Y$ 的概率分布为 $P\{Y=-1\}=\dfrac13$，$P\{Y=1\}=\dfrac23$，记 $Z=XY$.
\par（Ⅰ）求 $Z$ 的概率密度 $f_Z(z)$；
\par（Ⅱ）求 $\operatorname{Cov}(X,Z)$；
\par（Ⅲ）若 $Z_1,Z_2,\ldots,Z_n$ 为总体 $Z$ 的简单随机样本，求 $\lambda$ 的矩估计量 $\widehat\lambda$.
\end{question}
\WorkSpace
